\section{运行稳定性与结果核验}

\subsection{可复现运行}

问题一和问题二分别使用固定种子 0--5 的独立重启，问题三使用固定种子 0--1 并保留动态路线局部搜索。每次运行的配置、种子和结果写入结构化文件，最终只保留通过独立验证器的路线。图~\ref{fig:convergence} 展示真实运行中的多起点收敛记录，用于说明搜索过程和结果来源。

\begin{figure}[H]
\centering
\includegraphics[width=0.92\textwidth]{multistart_convergence.pdf}
\caption{固定随机种子多起点运行的真实收敛记录}
\label{fig:convergence}
\end{figure}

验证器从原始 Case 数据重新计算以下项目：每个 \texttt{task\_id} 是否恰好出现一次；合并连续相同物理点后的有效到达次数是否等于等级要求；路线是否从基地出发并返回；记录距离和时间是否与坐标、速度、服务时间重新计算结果一致；所有 $T_k$ 是否不超过 $9\ \mathrm{h}$。问题三进一步检查飞行线段和圆弧的几何合法性、禁飞时间重叠、巡检服务时段、基地停留和时间单调传播。

\subsection{理论下界与最低可行数量的区分}

理论下界仅由服务工作量和最远点往返给出，是搜索起点而非最终结论。由于本轮没有调用精确 MIP 不可行性证书，问题一和问题三的机队数量应表述为“最低已验证可行机队”。例如 Case1、Case3 和 Case4 的问题一理论工作下界为 2 架，而正式最好已验证结果分别需要 5、6 和 5 架；这反映了空间分布、任务不可相邻重复和基地闭环约束的共同影响。

\subsection{未运行实验的边界}

建模报告预留了速度、服务时间、禁飞半径和窗口平移等参数灵敏度实验，但本轮正式结果报告没有输出这些扰动实验。因此本文不把未实际运行的情景写成定量结论，仅报告已执行的固定种子多起点运行和正式路线验证。这样可以保持论文结论与真实结果文件的一致性。
